\section{Pemetaan Hamiltonian \& Hasil}

\begin{frame}
\centering
\highlight{\Huge Bagian 6: Pemetaan Hamiltonian \& Hasil}
\end{frame}

\subsection{Formulasi QUBO}
\begin{frame}{6.1 Formulasi QUBO berbasis Nash Equilibrium}
Fungsi potensial $\Phi(x)$ yang telah dianalisis melalui titik ekuilibrium (SBR) dimodifikasi menjadi sistem \textit{Quadratic Unconstrained Binary Optimization} (QUBO) dengan menambahkan penalti kardinalitas ($\lambda=5.0, K=2$).

\textbf{Fungsi Objektif QUBO:}
\begin{equation}
H_{QUBO}(x) = -\Phi(x) + \lambda \left( \sum_{i=1}^N x_i - K \right)^2
\end{equation}

\textbf{Elemen Matriks QUBO ($Q$):}
\begin{itemize}
    \item $Q_{ii} = -\tilde{\mu}_i + \frac{\gamma}{2} \tilde{\Sigma}_{ii} + \lambda (1 - 2K)$ (Linear)
    \item $Q_{ij} = \gamma \tilde{\Sigma}_{ij} + 2\lambda$ (Kuadratik)
\end{itemize}
$\Phi(x)$ yang bernilai minimum pada Nash Equilibrium murni kini menjadi target utama dalam sirkuit optimasi kuantum.
\end{frame}

\subsection{Hamiltonian Ising}
\begin{frame}{6.2 Parameter Ising & Representasi Matriks}
Transformasi $x_i \to \frac{I - Z_i}{2}$ memetakan masalah ke dalam Hamiltonian Ising yang terdiri dari kopling $J_{ij}$ dan bias $h_i$.

\textbf{Parameter Amplitudo Ising:}
\begin{equation}
J = 
\begin{bmatrix}
0.0 & 2.5 & 2.5 & 2.5 \\
2.5 & 0.0 & 2.5 & 2.5 \\
2.5 & 2.5 & 0.0 & 2.5 \\
2.5 & 2.5 & 2.5 & 0.0
\end{bmatrix}
, \quad h = 
\begin{bmatrix}
0,0020 \\ 0,0003 \\ 0,0015 \\ -0,0018
\end{bmatrix}
\end{equation}

\textbf{Strategi Quantum Warm-Start:}
Sirkuit VQE diinisialisasi menggunakan \textit{Nash Equilibrium} murni dari SBR ($0011$) untuk mempercepat konvergensi energi menuju \textit{ground state}.
\end{frame}

\subsection{Konvergensi & Hasil}
\begin{frame}{6.3 Hasil Optimasi & Konvergensi VQE}
Melalui algoritma SBR sebagai \textit{warm-start}, VQE berhasil menyempurnakan solusi optimasi dengan efisiensi yang lebih baik.

\textbf{Output Hasil:}
\begin{itemize}
    \item \textbf{Initial Seed (SBR):} $0011$ (ADRO \& TLKM)
    \item \textbf{Quantum Convergence:} VQE mengonfirmasi \textit{Ground State} dominan pada basis $|0011\rangle$ dengan energi $\approx -4,9989$.
\end{itemize}

\vspace{0.2cm}
\textbf{Signifikansi:} Integrasi pencarian ekuilibrium klasik ke dalam formulasi QUBO memastikan bahwa sistem kuantum memulai proses optimasi pada titik yang sudah mendekati solusi ideal.
\end{frame}

\begin{frame}
\centering \Huge
Terima Kasih
\end{frame}
